% Part: second-order-logic
% Chapter: metatheory
% Section: second-order-arithmetic

\documentclass[../../../include/open-logic-section]{subfiles}

\begin{document}

\olfileid{sol}{met}{spa}

\olsection{حساب الرتبة الثانية}

قد تقدم أن حساب بيانو~$\Th{PA}$ يشتمل على بديهيات~$\Th{Q}$
الثماني:
\begin{align*}
& \lforall[x][\eq/[x'][\Obj 0]]\\
& \lforall[x][\lforall[y][(\eq[x'][y'] \lif \eq[x][y])]]\\
& \lforall[x][(\eq[x][\Obj 0] \lor \lexists[y][\eq[x][y']])]\\ 
& \lforall[x][\eq[(x + \Obj 0)][x]]\\
& \lforall[x][\lforall[y][\eq[(x + y')][(x + y)']]]\\
& \lforall[x][\eq[(x \times \Obj 0)][\Obj 0]]\\
& \lforall[x][\lforall[y][\eq[(x \times y')][((x \times y) + x)]]]\\
& \lforall[x][\lforall[y][(x < y \liff \lexists[z][\eq[(z' + x)][y]])]]\\
\intertext{مضافًا إليها جميع الجمل من الصورة}
& (!A(\Obj 0) \land \lforall[x][(!A(x) \lif !A(x'))]) \lif \lforall[x][!A(x)].
\end{align*}
وليست الصورة الأخيرة جملة واحدة، بل «مخططًا» يولد من
!!{sentence}s لغة الحساب ما لا نهاية له: لكل !!{formula}~$!A(x)$
جملة تخصها. وهو ما نسميه \emph{مخطط بديهية الاستقراء} من
الرتبة الأولى. وأما حساب بيانو \emph{من الرتبة الثانية}~$\Th{PA^2}$
فيُجمع فيه الاستقراء في جملة واحدة؛ فتكون~$\Th{PA^2}$
مؤلفة من البديهيات الثماني المذكورة ومن \emph{بديهية الاستقراء}
من الرتبة الثانية:
\[
\lforall[X][((X(\Obj 0) \land \lforall[x][(X(x) \lif X(x'))]) \lif \lforall[x][X(x)])].
\]
ومعناها أن المجموعة الجزئية~$X$ من !!{domain} إذا احتوت
$\Assign{\Obj{0}}{M}$، وكانت متى احتوت~$x \in \Domain{M}$
احتوت خلفه~$\Assign{\prime}{M}(x)$، أي كانت «مغلقة تحت الخلف»،
لزم أن تحتوي كل ما في !!{domain}، أي~$X = \Domain{M})$.

وبهذه البديهية لا يكون في أي !!{structure} تحققها من
!!{element}s~$\Domain{M}$ إلا ما أوجبت البديهيات وجوده:
قيم~$\num{n}$ حيث~$n \in \Nat$.
فلا مكان للأعداد غير القياسية في نموذج لـ$\Th{PA^2}$.

\begin{thm}
\ollabel{thm:sol-pa-standard}
من~$\Sat{M}{\Th{PA^2}}$ يلزم~$\Domain{M} =
\Setabs{\Value{\num{n}}{M}}{n \in \Nat}$.
\end{thm}

\begin{proof}
ضع~$N = \Setabs{\Value{\num{n}}{M}}{n \in \Nat}$، ولنفترض
$\Sat{M}{\Th{PA^2}}$. لكل~$n \in \Nat$ يكون
$\Value{\num{n}}{M} \in \Domain{M}$؛ فثبت~$N \subseteq \Domain{M}$.

بقي إثبات الاحتواء المعاكس. نأخذ تعيين متغيرات~$s$
يحقق~$s(X) = N$. ومن الفرض،
\begin{align*}
\Sat{M}{\lforall[X][((X(\Obj 0) \land \lforall[x][(X(x)
      \lif X(x'))]) \lif \lforall[x][X(x)])]}, & \text{ومن ثم}\\
\Sat{M}{(X(\Obj 0) \land \lforall[x][(X(x) \lif X(x'))]) \lif
  \lforall[x][X(x)]}[s]. & 
\end{align*}
ولنثبت مقدم الشرطية. أوله ثابت لأن~$\Value{\Obj{0}}{M} \in N$،
فيكون~$\Sat{M}{X(\Obj{0})}[s]$.
وأما المقترن الثاني~$\lforall[x][(X(x) \lif X(x'))]$
فيُرضى كذلك: إذ لو كان~$x \in N$، لاقتضى تعريف~$N$
أن~$x = \Value{\num{n}}{M}$ لبعض~$n$.
فيكون~$\Assign{\prime}{M}(x) = \Value{\num{n+1}}{M} \in N$،
ومن ثم~$\Assign{\prime}{M}(x) \in N$.

فحصل~$\Sat{M}{X(\Obj 0) \land \lforall[x][(X(x) \lif X(x'))]}[s]$،
ولزم~$\Sat{M}{\lforall[x][X(x)]}[s]$.
ومعنى هذا أن~$x \in s(X) = N$ لكل~$x \in \Domain{M}$؛
فثبت~$\Domain{M} \subseteq N$.
\end{proof}

\begin{cor}
\ollabel{cor:sol-pa-categorical}
كل نموذجين لـ$\Th{PA^2}$ متماثلان بنيويًا.
\end{cor}

\begin{proof}
بمبرهنة~\olref{thm:sol-pa-standard} تستوعب
$\Value{\num{n}}{M}$ مجال كل نموذج لـ$\Th{PA^2}$.
وهو أيضًا نموذج لـ$\Th{Q}$؛ فينتج من
\olref[mod][mar][stm]{prop:thq-standard} أنه قياسي، أي متماثل
بنيويًا مع~$\Struct{N}$.
\end{proof}

جعلنا~$\Th{PA^2}$ مؤلفة من البديهيات الحسابية الثماني وبديهية
الاستقراء من الرتبة الثانية. غير أن ما في هذه الرتبة من قوة تعبيرية
يغنينا، في حساب بيانو من الرتبة الثانية، عن ست من تلك البديهيات؛
إذ تكفي البديهيتان \emph{الأوليان} مع الاستقراء.

\begin{prop}\ollabel{prop:sol-pa-definable}
لتكن~$\Th{PA^{2\dagger}}$ نظرية الرتبة الثانية المؤلفة من
البديهتين الحسابيتين الأوليين، وهما بديهيتا الخلف، ومن بديهية
الاستقراء من الرتبة الثانية. يمكن حينئذ تعريف~$\le$ و$+$ و$\times$
في~$\Th{PA^{2\dagger}}$.
\end{prop}

\begin{proof}
أما تعريف~$\le$ فيقتضي صيغة~$!A_\le(x, y)$ يكون بها
$\Sat{N}{!A_\le(\num n, \num m)}$ إذا وفقط إذا~$n \le m$.
فلنأخذ الصيغة
\[
!B(x, Y) \ident Y(x) \land \lforall[y][(Y(y) \lif Y(y'))]
\]
إن~$!B(\num n,Y)$ تعني أن~$Y$ تحتوي~$n$ ومغلقة تحت الخلف؛
فيلزم أن تحتوي كل~$m\ge n$. وإذا كان~$m<n$، حققت مجموعة الذيل
$\Setabs{k}{n\le k}$ الصيغة~$!B(\num n,Y)$ ولم تحتو~$m$.
ولذلك نضع~$!A_\le(x, y) \ident
\lforall[Y][(!B(x, Y) \lif Y(y))]$.
% تصحيح مصدر (0332-I1): استُبدل التكافؤ الزائف بالتوصيف اللازم وبمجموعة الذيل الشاهدة للاتجاه العكسي.

وأما الجمع، فإن~$k+l=m$ إذا وفقط إذا وجدت دالة~$u$
تحقق~$u(0)=k$ و$u(n')=u(n)'$ لكل~$n$، و$m = u(l)$.
فبهذا التكافؤ نعرّف الجمع في~$\Th{PA^{2\dagger}}$
بالصيغة:
\[
!A_+(x,y,z) \ident \lexists[u][(u(\Obj 0)=x \land \lforall[w][u(w')=u
 (w)'] \land u(y)=z)]
\] 
ومن البيان السابق يتضح أن~$\Sat{N}{!A_+(\num k, \num l, \num m)}$
إذا وفقط إذا كان~$k+l=m$.
\end{proof}

\begin{prob}
أتم ما بقي من برهان~\olref[sol][met][spa]{prop:sol-pa-definable}.
\end{prob}

\end{document}
